\documentclass[11pt]{article}
\usepackage[margin=1in]{geometry}
\usepackage{amsmath,amssymb,amsfonts,mathtools,bm}
\usepackage[T1]{fontenc}
\usepackage[utf8]{inputenc}
\title{Inline and Display Mathematics}
\author{Source-backed grouped formula sample}
\date{}
\begin{document}
\maketitle
\section{Expressions}
\paragraph{Expression 1.} The following expression is taken from a source-backed formula corpus.

\begin{align*}\underset{x^\prime \in \mathbb{R}^d}{\arg \max } \left[ \sum_{1\leq \ell \leq 2} \varphi_\ell(x^{\prime}, \lambda_\ell)\right] = \left(\sum_{1 \leq \ell \leq 2} \lambda_{\ell} V_{\ell, X X}^{-1}\right)^{-1} \left(\sum_{1 \leq \ell \leq 2} a_{\ell} V_{\ell, X X}^{-1} V_{\ell, X Y}\right),\end{align*}

\paragraph{Expression 2.} The following expression is taken from a source-backed formula corpus.

\begin{align*}{\bar G}^{(\zeta )}(\beta )={(\beta +i\pi )\sinh\left( {\displaystyle \pi\beta\over \displaystyle 2\pi +\zeta}\right) \over\Gamma \left( 1- {\displaystyle\pi-i\beta \over \displaystyle 2\pi+\zeta}\right)\Gamma \left( 1- {\displaystyle\pi+i\beta \over \displaystyle 2\pi+\zeta}\right)}\ ,\end{align*}

\paragraph{Expression 3.} The following expression is taken from a source-backed formula corpus.

\begin{align*}\underline D_\mu^\psi(\eta)\asymp_\times \min_{p\in P_{<\delta}} \liminf_{\xi\in G(p)} \left(\frac{1/\|\xi - \eta\|}{1/r_\xi} \psi_p(\|\xi - \eta\|) \right)^{k_p - \delta}= \min_{p\in P_{<\delta}} \left(\limsup_{\xi\in G(p)} \frac{r_\xi}{\theta_p(\|\xi - \eta\|)} \right)^{k_p - \delta}.\end{align*}

\paragraph{Expression 4.} The following expression is taken from a source-backed formula corpus.

\begin{align*} \sup_{\mathcal{T}>0}\ |\frac{1}{\mathcal{T}}\int_0^{\mathcal{T}} h(u(t),y_{\epsilon}(t),z_{\epsilon}(t))dt - \frac{1}{\mathcal{T}}\int_0^{\mathcal{T}}\tilde h(\mu(t),z(t))dt|\ \leq \alpha_h (\epsilon ), \ \ \ \ \ {\rm where} \ \ \ \ \ \lim_{\epsilon\rightarrow 0}\alpha_h (\epsilon) = 0\end{align*}

\paragraph{Expression 5.} The following expression is taken from a source-backed formula corpus.

\begin{align*}\gamma^{\mu\rho\nu}\nabla_{\rho}(\Psi^\mp_{\nu})^{A}-\gamma^{\mu\nu}\gamma^{5}\partial_{5}(\Psi^\pm_{\nu})^{A}-k\epsilon\tanh\left(k|y|-\frac{k\pi r_{c}}{2}\right)\gamma^{\mu\nu}\gamma^{\hat{5}}(\Psi^\pm_{\nu})^{A}-\frac{3}{2}k(\sigma_1)^{A}_{\;B}\gamma^{\mu\nu}(\Psi^\mp_{\nu})^{B}=0\ .\end{align*}
\end{document}
